\section{De características manuales a auto-supervisión: ¿qué es lo que realmente aprenden los modelos?}

Esta clase de apertura del curso V6 no simplifica el Transformer reduciéndolo únicamente a una fórmula de atención. Los dos ponentes primero rastrean cómo la señal de aprendizaje evolucionó desde características manuales y etiquetas supervisadas hacia objetivos de auto-supervisión, y luego conectan datos, post-entrenamiento, agentes, multimodalidad, memoria y alineación en un mismo mapa del sistema. Al leer estas notas, debe mantenerse preguntando constantemente tres cuestiones: ¿de dónde proviene la información?, ¿qué recompensan los gradientes? y, ¿qué mecanismos externos al modelo añaden los sistemas una vez desplegados?

\subsection{Promesa del curso: mecanismo, aplicación y límites deben examinarse conjuntamente}

Los puntos clave iniciales sitúan la atención, la escalabilidad, las aplicaciones interdominio, los ponentes de vanguardia y los problemas abiertos en la misma página. Este orden define el principio organizador de las notas: primero explicar el mecanismo, luego contrastar con la evidencia y finalmente explicitar los límites. Conocer solo la forma tensorial de un bloque es insuficiente para explicar por qué un modelo tiene éxito sobre un determinado conjunto de datos; ver únicamente una demostración exitosa tampoco basta para juzgar si es fiable en nuevas distribuciones, contextos largos o flujos de trabajo reales.

\lecturefigure{slide-011.jpg}{Objetivos del curso: desde el mecanismo del Transformer hasta aplicaciones de vanguardia, limitaciones y problemas abiertos}{CS25 V6 oficial deck, p.~11}

\teachervoice{00:06:12--00:07:10，el docente enfatiza que el curso no se limita a explicar la atención, sino que busca que los estudiantes accedan a aplicaciones de vanguardia, limitaciones clave y problemas abiertos en los que aún pueden contribuir. Nota de clase: cada método posterior debe responder a "¿qué capa del sistema ha cambiado?".}

\begin{importantbox}{Las seis capas de interfaz de los sistemas Transformer modernos}
\begin{enumerate}
  \item \textbf{Representación}: cómo la entrada original se convierte en tokens y vectores;
  \item \textbf{Arquitectura}: cómo fluye la información entre secuencias, capas y modalidades;
  \item \textbf{Pre-entrenamiento}: qué capacidades por defecto moldean la distribución de datos y la función objetivo;
  \item \textbf{Post-entrenamiento}: cómo la supervisión, las preferencias, la búsqueda y el verificador modifican el comportamiento;
  \item \textbf{Bucle del sistema}: cómo la recuperación, las herramientas, la memoria y la retroalimentación ambiental ingresan al ciclo cerrado;
  \item \textbf{Gobernanza/evidencia}: cómo se definen la evaluación, la explicabilidad, la alineación y los límites de despliegue.
\end{enumerate}
Un nuevo término que no pueda aterrizar en alguna de estas interfaces generalmente no ha sido verdaderamente explicado.
\end{importantbox}

\subsection{Tres transformaciones: ingeniería de características, representación supervisada y representación auto-supervisada}

El aprendizaje automático temprano dependía de que los investigadores diseñaran manualmente las características, para luego que un modelo relativamente superficial las mapeara a etiquetas. El aprendizaje profundo delegó el aprendizaje de la representación a la red, pero aún exigía abundantes etiquetas manuales. El aprendizaje auto-supervisado continuó desplazando la fuente de supervisión: las etiquetas ya no provienen de categorías humanas, sino que se construyen a partir de la entrada misma, como ocultar tokens, predecir el siguiente paso, restaurar imágenes perturbadas o distinguir muestras coincidentes de no coincidentes. La diferencia clave entre los tres paradigmas no es "qué tan profunda es la red", sino \emph{quién define la tarea de aprendizaje y cuánta estructura original puede ver el modelo}.

\lecturefigure{slide-013.jpg}{Aprendizaje automático temprano: características manuales, optimizador y costosas etiquetas manuales}{CS25 V6 oficial deck, p.~13}

\lecturefigure{slide-014.jpg}{Profundización supervisada: la red comienza a asumir el aprendizaje de representación, pero sigue dependiendo de etiquetas}{CS25 V6 oficial deck, p.~14}

De las diapositivas 13 a la 16 se debe leer según el cuello de botella de información: las características manuales comprimen los datos originales en dimensiones que el investigador considera útiles; la red supervisada amplía el alcance de la representación aprendible; y el entrenamiento end-to-end reduce aún más las interfaces intermedias humanas. El costo también cambia: aumenta la libertad del modelo, lo que requiere más datos y cómputo; si la distribución de entrenamiento presenta sesgos, la red aprenderá esos sesgos junto con los patrones válidos.

\lecturefigure{slide-015.jpg}{Estado transicional del aprendizaje supervisado: los datos originales aún pasan primero por una representación manual}{CS25 V6 oficial deck, p.~15}

\lecturefigure{slide-016.jpg}{Aprendizaje supervisado end-to-end: la entrada original ingresa directamente a la red profunda}{CS25 V6 oficial deck, p.~16}

\begin{knowledgebox}{Lectura de la figura: por qué valen las cuatro páginas progresivas}
Primero compare quién construye la representación, luego desde dónde proviene la señal de supervisión. En la diapositiva 13, los humanos definen las características; en las diapositivas 14--15 coexisten la red y la interfaz manual; en la diapositiva 16, la red consume directamente los datos brutos. La figura respalda el desplazamiento interno gradual de la "interfaz de aprendizaje de representación", pero no prueba que el conocimiento humano carezca de valor. Los sistemas modernos aún inyectan abundante diseño humano mediante tokenizadores, filtrado de datos, priorización arquitectónica, pérdida y APIs de herramientas.
\end{knowledgebox}

El objetivo del aprendizaje auto-supervisado puede abstraerse como construir una parte predecible $y=g(x)$ y un contexto visible $\tilde{x}=c(x)$ a partir de una muestra $x$, para luego optimizar
\[
\mathcal{L}_{\mathrm{SSL}}(\theta)
=-\log p_{\theta}\!\left(y\mid \tilde{x}\right).
\]
$g$ representa la regla para construir el objetivo, como el siguiente token o un parche oculto; $c$ representa el contexto conservado para el modelo; $\theta$ son los parámetros. Aunque la supervisión "proviene de los propios datos", $g$ y $c$ aún deciden qué información se incentiva que el modelo conserve.

\lecturefigure{slide-017.jpg}{Aprendizaje auto-supervisado: construcción del objetivo de entrenamiento desde el interior de los datos originales}{CS25 V6 oficial deck, p.~17}

\teachervoice{00:07:10--00:08:13，los ponentes resumen la línea histórica como características diseñadas a mano, aprendizaje de representación supervisado y auto-supervisión. El docente enfatiza que el aprendizaje auto-supervisado reduce las etiquetas manuales, pero no elimina el sesgo inductivo traído por el objetivo.}

\begin{warningbox}{"Sin etiquetas manuales" no equivale a "sin suposiciones humanas"}
La proporción de ocultamiento, la dirección de predicción, las muestras negativas, el aumento de datos y la longitud de la ventana alteran la señal aprendible. Si el aumento destruye factores relevantes para la tarea, un modelo contrastivo será recompensado por ignorar información verdaderamente útil; si los datos del siguiente token están llenos de plantillas y repeticiones, el modelo priorizará aprender formas de alta frecuencia. El aprendizaje auto-supervisado es una forma de construir etiquetas, no una ley natural neutral al valor.
\end{warningbox}

Desde la perspectiva del cuello de botella de información, el objetivo decide qué información merece ser conservada. La clasificación de sentimientos solo requiere conservar las direcciones relacionadas con la etiqueta, y puede descartar activamente entidades, secuencia temporal y detalles gramaticales; el objetivo de reconstrucción necesita conservar más información local, pero podría codificar también una gran cantidad de ruido de píxeles; la predicción del siguiente token debe comprimir las leyes estadísticas útiles para futuros tokens, pero no recibe recompensa directa por estados del mundo real que no cambian la probabilidad del token. Por tanto, al evaluar una representación no basta mirar solo una puntuación downstream, sino que se deben revisar sondas lineales, transferencia con pocos ejemplos (\textit{few-shot}), cambio de distribución e intervenibilidad. Una representación que funcione bien en clasificación no significa que sea adecuada para planificación; un modelo generativo que pueda reconstruir detalles no implica que su latente ya haya abstraído objetos y causalidad.

\subsection{Casos lingüísticos y visuales: el objetivo de predicción determina la granularidad de la representación}

El caso lingüístico pasa de la clasificación de sentimientos a la predicción del siguiente token. Lo primero solo exige comprimir un texto en pocas etiquetas; lo segundo requiere predecir la distribución condicional en cada posición, por lo que la supervisión es más densa. Sea $x_{1:T}$ la secuencia de tokens; el modelo de lenguaje auto-regresivo optimiza
\[
\mathcal{L}_{\mathrm{NTP}}(\theta)
=-\sum_{t=1}^{T}\log p_{\theta}(x_t\mid x_{<t}).
\]
$x_t$ es el $t$-ésimo token, $x_{<t}$ es el contexto histórico. Cada posición contribuye a la pérdida, pero el objetivo sigue requiriendo solo predecir la secuencia observada, sin exigir directamente que el modelo proporcione fuentes factuales, ejecute herramientas o mantenga memoria a largo plazo.

\lecturefigure{slide-018.jpg}{Caso lingüístico: paso de etiquetas esparsas de sentimiento a predicción densa de secuencias}{CS25 V6 official deck, p.~18}

\lecturefigure{slide-020.jpg}{Estado completo del auto-supervisión lingüística: predecir el siguiente token bajo condición del contexto}{CS25 V6 official deck, p.~20}

\teachervoice{00:08:13--00:09:57，el docente contrasta la clasificación de sentimientos con el modelado de lenguaje: la clasificación proporciona solo una etiqueta de baja banda ancha, mientras que la predicción de secuencias genera supervisión en cada posición. Experiencia práctica: una supervisión más densa no implica que el objetivo cubra automáticamente la factualidad y la causalidad.}

El auto-supervisión visual suele elegir entre reconstrucción y aprendizaje contrastivo. El objetivo de reconstrucción exige que el decodificador restituya la entrada, lo que facilita conservar detalles de píxeles; el objetivo contrastivo acerca las vistas aumentadas del mismo objeto y aleja muestras distintas, enfatizando más la representación invariante. Una pérdida InfoNCE simplificada es
\[
\mathcal{L}_{\mathrm{NCE}}
=-\log\frac{\exp(\operatorname{sim}(z_i,z_i^+)/\tau)}
{\sum_j \exp(\operatorname{sim}(z_i,z_j)/\tau)}.
\]
$z_i$ es la representación ancla, $z_i^+$ la muestra positiva y $z_j$ las muestras candidatas; $\tau$ es la temperatura. Lo clave no es memorizar la fórmula, sino comprender que la definición de muestras positivas y negativas decide qué cambios se exige al modelo ignorar.

\lecturefigure{slide-021.jpg}{Caso visual: dos interfaces de supervisión entre reconstrucción con autoencoder y aprendizaje contrastivo}{CS25 V6 official deck, p.~21}

\lecturefigure{slide-022.jpg}{Autoencoder enmascarado: predecir contenido visual faltante tras ocultar parches}{CS25 V6 official deck, p.~22}

\begin{knowledgebox}{Lectura de la figura: límites de evidencia entre reconstrucción y contraste}
Primero observe el espacio de salida: los modelos de reconstrucción predicen detalles de entrada, mientras que los modelos contrastivos predicen relaciones entre muestras. Luego examine la invariancia: el segundo comprimirá activamente los cambios considerados "poco importantes" en el aumento. La figura respalda las diferencias entre dos objetivos de representación, pero no permite afirmar solo con un diagrama cual de ambos es generalmente mejor; tareas downstream, escala de datos, aumentos y capacidad del decodificador alterarán la conclusión.
\end{knowledgebox}

\subsection{Resumen de la sección}

La evolución de los paradigmas de aprendizaje se puede comprimir en "el desplazamiento interno continuo de la interfaz de supervisión": primero los humanos crean características, luego las redes aprenden representaciones y finalmente los propios datos generan objetivos densos. Pero cualquier objetivo solo recompensa patrones estadísticos específicos. Antes de entender el Transformer, hay que comprender que recibe tokens, optimiza predicciones condicionales y no obtiene automáticamente hechos, acciones ni coherencia a largo plazo.


